\section{Conclusion and Future Work}\label{sec:conclusion}
This project evaluates conventional machine learning for hourly total electricity load in Tetouan using \num{52416} ten-minute observations aggregated into \num{8736} hourly means. Its contribution is a documented experimental pipeline combining seasonal references, chronological partitions, five expanding development folds, randomized parameter search, residual-target comparisons, feature ablation and evaluation diagnostics.

Absolute-target SVR minimizes development mean CV RMSE in both implemented settings. On the separate December period, the selected h24 model achieves RMSE \num{1624.6} kW and MAPE \num{1.78}\%, reducing RMSE by \num{50.71}\% against Naive-24h. The selected h1 model reaches RMSE \num{1195.8} kW and MAPE \num{1.41}\%. Residual SVR has lower observed h24 test RMSE, \num{1419.1} kW, but remains a test-best observation rather than a replacement for the CV-selected model.

The experiments show why feature and evaluation choices matter alongside algorithm choice. Daily lags provide the largest tested sequential feature improvement for RF/XGB; adding weekly and rolling summaries can worsen their scores under fixed D-selected parameters. Residual targets help most tuned families on development, while temporal folds reveal summer weaknesses missed by a single November split. S1--S4 demonstrate sensitivity to task and protocol, but their simultaneous changes prevent attributing the entire score difference to leakage.

These findings are conditional on observed target-hour weather and rapidly available historical loads. One year, one December holdout, reused development folds, descriptive test exposure and documented test reruns limit generalization claims. The report consequently supports accurate reporting under explicit information assumptions, rather than superiority over heterogeneous published results or deployment readiness.

Future work should first replace oracle weather with forecasts archived at issuance and incorporate measured load-delivery delays. A further independent year or several rolling seasonal holdouts would allow evaluation of summer peaks and holiday behavior after freezing development choices. Nested temporal selection or an additional outer validation layer could separate tuning from performance estimation. Retuned weather-free and weekly-feature ablations, with matched samples and availability, would better assess feature utility.

Additional horizons and explicitly stated issuance-to-target leads would clarify the relation between forecast difficulty and available history. Candidate extensions include classical time-series and alternative ML models, event-specific calendar encodings, per-zone forecasting, uncertainty intervals and explanations for the selected SVR. A later deployment study should assess monitoring, drift, fallback rules and computational cost. These are proposed investigations; none is presented as an implemented result of the current repository.
